\chapter{Phương pháp đề xuất}
\section{Tổng quan phương pháp}
Hệ thống được chia thành các module độc lập nhưng nối tiếp nhau. Module 1 chuẩn bị dữ liệu. Module 2 huấn luyện sentiment model. Module 3 phân tích thống kê review. Module 4 xây dựng retrieval và RAG answer. Module 5 tổng hợp kết quả và phân tích lỗi. Module 6 là hướng nâng cấp retrieval.

\section{Module 1 - Data Layer}
Module 1 đọc dữ liệu review, làm sạch text, chuẩn hóa rating, xử lý duplicate và tạo nhãn sentiment. File chính là \code{src/module1_data_layer.py}. File này đóng vai trò nền móng vì dữ liệu sai sẽ kéo theo model sai và RAG retrieval kém.

Các function chính trong Module 1 gồm nhóm đọc dữ liệu, nhóm chuẩn hóa text/rating, nhóm ánh xạ rating sang sentiment, nhóm xử lý duplicate/conflict duplicate và nhóm chia train/valid/test. Output quan trọng gồm \code{data/processed/sentiment_reviews.csv}, \code{data/processed/sentiment_reviews_5class.csv}, \code{data/processed/rag_corpus_module4.csv} và \code{data/reports/data_summary.json}. Về mặt phương pháp, Module 1 biến dữ liệu thô thành dữ liệu sạch có thể huấn luyện và truy xuất.

\section{Module 2 - Sentiment Classification}
Module 2 huấn luyện các mô hình phân loại cảm xúc. File chính là \code{src/module2_baseline_sentiment_classification.py}. Module này thử ba hướng: majority baseline, TF-IDF + Logistic Regression và TF-IDF + Linear SVM. Mô hình được chọn theo validation macro-F1.

Majority baseline cho biết mốc tối thiểu. Logistic Regression là baseline tuyến tính có xác suất. Linear SVM là baseline mạnh cho TF-IDF sparse vector. Các artifact chính là \code{models/module2/best_model.joblib}, \code{models/module2/tfidf_linearsvm.joblib}, \code{models/module2/tfidf_logreg.joblib} và \code{models/module2/label_mapping.json}.

Khi inference, web backend load \code{best_model.joblib}. Vì best model là pipeline gồm TF-IDF vectorizer và classifier, người dùng chỉ cần nhập text review, hệ thống trả về sentiment. Nếu model có decision score hoặc margin, backend có thể hiển thị thêm mức tự tin tương đối.
\subsection{Hyperparameter thật trong Module 2}
File \filepath{src/module2_baseline_sentiment_classification.py} dùng random seed 42 và chọn mô hình theo validation macro-F1. Bảng \ref{tab:module2-hyperparams} tóm tắt cấu hình quan trọng lấy từ code và kết quả tuning.
\begin{table}[H]
\centering
\caption{Cấu hình/hyperparameter chính của Module 2}
\label{tab:module2-hyperparams}
\small
\begin{tabularx}{\textwidth}{L{0.28\textwidth}Y}
\toprule
Thành phần & Cấu hình \\
\midrule
TF-IDF search space & \code{ngram_range}: (1,1) hoặc (1,2); \code{min_df}: 1, 2, 3; \code{max_df}: 0.95; \code{sublinear_tf}: True \\
Logistic Regression & \code{solver=lbfgs}, \code{max_iter=2000}, \code{C}: 0.5/1.0/2.0, \code{class_weight}: None hoặc balanced, \code{random_state=42} \\
Linear SVM & \code{LinearSVC}, \code{C}: 0.5/1.0/2.0, \code{class_weight}: None hoặc balanced, \code{random_state=42} \\
Best Logistic Regression & \code{C=2.0}, \code{class_weight=balanced}, \code{min_df=1}, \code{ngram_range=(1,2)} \\
Best Linear SVM & \code{C=1.0}, \code{class_weight=None}, \code{min_df=1}, \code{ngram_range=(1,2)} \\
Tiêu chí chọn model & Validation macro-F1, sau đó đánh giá cuối trên test \\
\bottomrule
\end{tabularx}
\end{table}

\subsection{Pseudocode huấn luyện sentiment}
\begin{lstlisting}[language=Python]
train, valid, test = load_sentiment_splits()
for vectorizer_params in tfidf_grid:
    X_train = fit_tfidf(train.text, vectorizer_params)
    for clf_params in model_grid:
        model = train_classifier(X_train, train.label, clf_params)
        valid_pred = model.predict(valid.text)
        score = macro_f1(valid.label, valid_pred)
        keep_best_model_if_score_is_higher(model, score)
test_pred = best_model.predict(test.text)
save_model_and_metrics(best_model, test_pred)
\end{lstlisting}


\section{Module 3 - Review Analytics}
Module 3 không huấn luyện model mới. Nó dùng dữ liệu sentiment để tạo thống kê và biểu đồ. Các output gồm \code{overall_sentiment_distribution.csv}, \code{rating_sentiment_crosstab.csv}, \code{top_negative_reviews.csv} và các hình trong \code{figures/module3/}.

Vai trò của Module 3 là giúp người đọc hiểu dữ liệu. Ví dụ, nếu positive và negative khá cân bằng nhưng neutral ít hơn, ta có thể dự đoán neutral là lớp khó. Crosstab rating-sentiment giúp kiểm tra ánh xạ rating sang sentiment. Top negative reviews giúp minh họa dữ liệu thực tế.

\section{Module 4 - Retrieval/RAG}
Module 4 là phần hỏi đáp dựa trên review. Các file chính gồm \code{src/module4_prepare_inputs.py}, \code{src/module4_build_index.py}, \code{src/module4_retrieve.py}, \code{src/module4_rag_pipeline.py} và \code{src/module4_evaluate.py}. Phiên bản giao diện thử nghiệm dùng thêm file nâng cấp \code{src/module4_retrieve_upgraded.py} và \code{src/module4_rag_pipeline_upgraded.py}.

\subsection{Chuẩn bị input}
\code{module4_prepare_inputs.py} chuẩn hóa corpus cho retrieval. Nó đảm bảo mỗi document có text truy xuất, metadata như product name, category, rating, sentiment và doc id. Output là corpus module4 để build index.

\subsection{Build index}
\code{module4_build_index.py} dùng multilingual-E5 để encode từng review. Document được thêm prefix \code{passage: }. Embedding được normalize và đưa vào FAISS \code{IndexFlatIP}. Output là \code{review_faiss.index}, \code{review_metadata.parquet} và \code{module4_index_config.json}.

\subsection{Retrieve và RAG pipeline}
\code{module4_retrieve.py} và bản nâng cấp \code{module4_retrieve_upgraded.py} load FAISS index, metadata và embedding model. Query được thêm prefix \code{query: }, encode thành vector, rồi tìm top-k review gần nhất. Bản nâng cấp thêm query expansion tiếng Việt, metadata filtering theo sentiment/category/rating và lightweight rerank theo từ khóa/aspect.

\code{module4_rag_pipeline.py} và \code{module4_rag_pipeline_upgraded.py} tạo câu trả lời từ retrieved reviews. Câu trả lời có cấu trúc: tóm tắt, vấn đề chính, bằng chứng tiêu biểu, độ tin cậy và lưu ý không suy diễn ngoài dữ liệu. Vì generator là template-based, hệ thống không cần API key và ít rủi ro hallucination hơn LLM tự do.


\subsection{Pseudocode retrieval/RAG}
\begin{lstlisting}[language=Python]
query = normalize(user_query)
expanded_query, inferred_filters = infer_filters_and_expand_query(query)
q_vec = e5.encode("query: " + expanded_query, normalize=True)
scores, ids = faiss_index.search(q_vec, candidate_k)
candidates = metadata.iloc[ids]
candidates = apply_metadata_filters(candidates, inferred_filters)
candidates = lightweight_rerank(candidates, query, aspect_keywords)
evidence = candidates.head(top_k)
answer = template_generate_answer(query, evidence, citations=True)
\end{lstlisting}
Điểm cần nhấn mạnh là bước generate không gọi LLM bên ngoài. Nó chỉ tổng hợp nội dung từ evidence đã truy xuất và gắn citation tương ứng.

\section{Module 5 - Evaluation \& Report}
Module 5 tổng hợp kết quả từ các module trước. Các file \code{src/module5_collect_results.py}, \code{src/module5_manual_retrieval_eval.py}, \code{src/module5_error_analysis.py}, \code{src/module5_generate_figures.py} và \code{src/module5_generate_report_sections.py} phục vụ thu thập metric, tạo bảng/hình, phân tích lỗi và sinh nội dung báo cáo.

Manual Precision@k là cách đánh giá relevance: với mỗi query, từng review trong top-k được gán nhãn liên quan hoặc không liên quan. File \code{manual_precision_at_k_sheet.csv} là bảng ghi nhãn relevance, còn \code{manual_precision_overall.csv} là summary. Cách đánh giá này tốn công nhưng cần thiết vì retrieval relevance không thể chỉ đo bằng score embedding.

\section{Module 6 - Aspect-aware Hybrid Retrieval}
Module 6 bổ sung BM25, Dense FAISS, RRF fusion, optional reranker và rule-based aspect sentiment. File chính gồm \code{src/module6_hybrid_retrieval.py}, \code{src/module6_reranker.py}, \code{src/module6_aspect_sentiment.py} và \code{src/module6_query_intent.py}.

Lý do cần Module 6 là dense retrieval có thể bỏ sót các cụm complaint ngắn như “sai size”, “móp hộp”, “shop không trả lời”, “giao sai”. BM25 bắt keyword tốt hơn; dense retrieval bắt ngữ nghĩa tốt hơn. RRF fusion kết hợp cả hai. Aspect rules giúp hệ thống hiểu query theo khía cạnh như \code{chat_luong}, \code{giao_hang}, \code{dong_goi}, \code{gia}, \code{dich_vu_shop}, \code{mau_ma_size_mau}. Cần nhấn mạnh Module 6 là MVP/experimental enhancement, không phải supervised ABSA hoàn chỉnh và không phải chatbot LLM.

\section{Bảng artifact chính}
\begin{longtable}{L{0.30\textwidth}L{0.60\textwidth}}
\caption{Artifact chính của từng module}\label{tab:artifacts}\\
\toprule
Artifact & Ý nghĩa \\
\midrule
\code{best_model.joblib} & Model sentiment chính được chọn theo validation macro-F1. \\
\code{tfidf_linearsvm.joblib} & Pipeline TF-IDF + Linear SVM. \\
\code{tfidf_logreg.joblib} & Pipeline TF-IDF + Logistic Regression. \\
\code{label_mapping.json} & Mapping nhãn sentiment. \\
\code{review_faiss.index} & Vector index dùng cho retrieval. \\
\code{review_metadata.parquet} & Metadata review tương ứng với vector trong FAISS. \\
\code{module4_index_config.json} & Cấu hình embedding, index, số documents. \\
\code{manual_precision_overall.csv} & Kết quả Precision@k thủ công của retrieval baseline. \\
\code{last_hybrid_rerank_demo.json} & Output thử nghiệm mới nhất của Module 6. \\
\bottomrule
\end{longtable}

Đường dẫn đầy đủ của các artifact được trình bày trong phụ lục để tránh làm bảng chính bị quá tải bởi các chuỗi path dài.

\section{Luồng inference trong giao diện thử nghiệm}
Khi người dùng nhập review ở tab Sentiment, frontend gửi request đến \code{/api/sentiment}. Backend load \code{best_model.joblib}, chạy pipeline TF-IDF + classifier và trả nhãn. Đây là inference một văn bản độc lập.

Khi người dùng nhập query ở tab RAG, frontend gửi request đến \code{/api/rag}. Backend không dùng sentiment classifier để phân loại query như review. Thay vào đó, backend dùng intent/aspect rules để hiểu query: nếu có “phàn nàn”, “lỗi”, “hỏng”, “chậm” thì ưu tiên sentiment negative; nếu có “giao hàng”, “ship”, “vận chuyển” thì aspect là giao hàng; nếu có “size”, “rộng”, “chật” thì có thể ưu tiên Fashion. Sau đó query được mở rộng, encode bằng E5, search FAISS, lọc/rerank và tạo answer.

\section{Kiểm soát chất lượng câu trả lời RAG}
RAG answer trong đề tài tuân theo bốn phần: tóm tắt, vấn đề chính, bằng chứng tiêu biểu và độ tin cậy. Template không được bịa thêm thông tin ngoài review. Nếu top-k review không liên quan rõ, answer phải nói “Chưa đủ bằng chứng rõ ràng”. Mỗi ý chính chỉ nên có 2--3 citation mạnh để tránh rối mắt.

Độ tin cậy trong UI không phải xác suất thống kê tuyệt đối. Nó là đánh giá heuristic dựa trên số lượng evidence, mức độ khớp query/aspect và sentiment phù hợp. Vì vậy, độ tin cậy cần được diễn giải là mức tin cậy của câu trả lời dựa trên evidence retrieved, không phải metric model giống accuracy.

\section{Vì sao Module 6 chưa thay thế hoàn toàn Module 4}
Module 4 là pipeline chính ổn định cho giao diện thử nghiệm. Module 6 là hướng nâng cấp kỹ thuật. Lý do không thay thế hoàn toàn là Module 6 cần thêm benchmark lớn hơn và kiểm thử latency. Hybrid retrieval có thể tốt hơn trong complaint query, nhưng cũng phức tạp hơn: cần BM25 index, dense retriever, fusion, reranker và aspect rules. Trong báo cáo, Module 6 được trình bày như phần mở rộng chứng minh hướng cải tiến, không phải khẳng định final production retriever.

\section{Tính tái lập}
Hệ thống lưu artifact và config để có thể tái lập. \code{module4_index_config.json} lưu embedding model, index type, dimension, số documents và path. \code{metrics_summary.csv} lưu metric sentiment. \code{manual_precision_overall.csv} lưu kết quả retrieval manual. Việc lưu file trung gian giúp báo cáo không phụ thuộc vào việc chạy lại toàn bộ pipeline, đồng thời vẫn có thể kiểm tra nguồn số liệu.

\section{Mô tả chi tiết các file code chính}
\subsection{File \code{src/module1_data_layer.py}}
File này làm nhiệm vụ chuẩn hóa dữ liệu đầu vào. Về mặt pipeline, nó nằm trước tất cả các module khác. Input là dữ liệu review thô và các trường rating/comment. Output là các file processed và summary. Các nhóm hàm quan trọng thường gồm: hàm đọc dữ liệu, hàm chuẩn hóa text, hàm chuẩn hóa rating, hàm tạo sentiment label, hàm xử lý duplicate và hàm chia dữ liệu. Nếu file này sai, mọi metric phía sau đều có nguy cơ sai vì model học từ dữ liệu không sạch.

\subsection{File \code{src/module2_baseline_sentiment_classification.py}}
File này huấn luyện và đánh giá các baseline sentiment. Input là tập train/valid/test từ Module 1. Output là model artifact và metric trong \code{results/module2/}. Về mặt học thuật, file này thể hiện quy trình Machine Learning rõ ràng: định nghĩa baseline, huấn luyện, chọn model, đánh giá trên test và lưu model.

\subsection{File \code{src/module4_prepare_inputs.py}}
File này chuyển dữ liệu review sang dạng phù hợp cho retrieval. Nó cần tạo trường text truy xuất đầy đủ nhưng không quá nhiễu. Một document RAG tốt nên có thông tin domain/category, product name, rating và comment. Nếu chỉ embed comment quá ngắn, retrieval có thể thiếu ngữ cảnh. Nếu embed quá nhiều metadata không liên quan, retrieval có thể bị nhiễu.

\subsection{File \code{src/module4_build_index.py}}
File này là bước build artifact nặng nhất của RAG. Nó load embedding model, encode từng passage, normalize embedding và ghi FAISS index. Bước này chỉ cần chạy lại khi corpus, embedding model hoặc cách tạo retrieval text thay đổi. Nếu chỉ sửa frontend hoặc format answer, không cần build lại index.

\subsection{File \code{src/module4_retrieve_upgraded.py}}
File nâng cấp này phục vụ phiên bản giao diện thử nghiệm. Nó giữ logic dense retrieval nhưng thêm query expansion, filter, loại category nhiễu nếu cần và rerank nhẹ theo aspect. Đây là nơi query người dùng được chuyển thành retrieval action.

\subsection{File \code{src/module4_rag_pipeline_upgraded.py}}
File này tạo answer từ retrieved reviews. Nó không gọi LLM. Nó format câu trả lời thành các phần: tóm tắt, vấn đề chính, bằng chứng, độ tin cậy. Đây là file quan trọng để chứng minh RAG của hệ thống không bịa ngoài dữ liệu.

\subsection{File \code{backend/app.py}}
File backend kết nối các artifact với giao diện. Nó load model sentiment, load retriever/RAG pipeline, định nghĩa API và trả JSON. Khi API lỗi, file này và log server là hai nguồn kiểm tra chính.

\subsection{File \code{frontend/app.js}}
File frontend quản lý state UI. Nó đảm bảo input sentiment riêng với input RAG, lưu last-run query, hiển thị cảnh báo khi query thay đổi, gọi API và render answer/evidence. Đây là phần quan trọng để hệ thống không gây hiểu nhầm giữa review cần phân loại và câu hỏi RAG.

\section{Vai trò kết hợp giữa mô hình phân loại cảm xúc và tầng RAG}
Nếu chỉ có sentiment model, hệ thống chỉ trả lời được “review này tiêu cực hay tích cực”. Nếu chỉ có RAG, hệ thống có thể tìm review nhưng thiếu lớp phân tích cảm xúc có giám sát. Đề tài yêu cầu ứng dụng NLP và RAG trong phân tích cảm xúc, nên hai phần cần liên kết: sentiment model tạo năng lực phân loại review, còn RAG dùng corpus review và metadata sentiment để trả lời câu hỏi tổng hợp có bằng chứng.

Trong giao diện thử nghiệm, sentiment được đặt trước vì đó là bài toán chính của môn NLP. RAG xuất hiện sau như lớp hỏi đáp và giải thích. Cách trình bày này phù hợp với tên đề tài: phân tích cảm xúc là trọng tâm, retrieval-augmented generation giúp mở rộng sang insight có citation.


\subsection{Pseudocode Module 6 hybrid retrieval}
\begin{lstlisting}[language=Python]
bm25_results = bm25_search(query, top_n)
dense_results = faiss_dense_search(query, top_n)
fused = reciprocal_rank_fusion([bm25_results, dense_results], k=60)
aspect = detect_aspect(query)
fused = apply_aspect_rules(fused, aspect)
if cross_encoder_available:
    final = cross_encoder_rerank(query, fused)
else:
    final = lexical_lightweight_rerank(query, fused)
return final.head(top_k)
\end{lstlisting}
Module 6 vì vậy là một bản nâng cấp retrieval có kiểm soát, không phải chatbot LLM hoàn chỉnh và cũng chưa phải ABSA supervised.
